\usepackage[T1]{fontenc}
\usepackage[utf8]{inputenc}
\usepackage{lmodern}
\usepackage[a4paper,margin=2.3cm,marginparwidth=1.6cm]{geometry}
\usepackage{amsmath,amssymb,amsthm,mathtools}
\usepackage{bm}
\usepackage{cancel}
\usepackage{enumitem}
\usepackage{xcolor}
\usepackage{tikz}
\usetikzlibrary{arrows.meta,calc,positioning,decorations.pathreplacing,patterns,angles,quotes}
\usepackage{pgfplots}
\pgfplotsset{compat=1.18}
\usepackage[most]{tcolorbox}
\usepackage{array,booktabs,multirow}
\usepackage{hyperref}
\hypersetup{colorlinks=true,linkcolor=blue!50!black}

\usepackage{titlesec}
\usepackage{marginnote}
\usepackage{needspace}
\titleformat{\section}{\needspace{8\baselineskip}\Large\bfseries\color{blue!45!black}}{\thesection}{0.8em}{}[{\color{blue!45!black}\titlerule[1pt]}]
\titlespacing*{\section}{0pt}{2.2em}{1.2em}
\titleformat{\subsection}{\needspace{6\baselineskip}\large\bfseries\color{blue!35!black}}{\thesubsection}{0.7em}{}[{\color{gray!60}\titlerule[0.5pt]}]
\titleformat{name=\subsection,numberless}{\needspace{6\baselineskip}\large\bfseries\color{blue!35!black}}{}{0pt}{}[{\color{gray!60}\titlerule[0.5pt]}]
\titlespacing*{\subsection}{0pt}{1.6em}{0.8em}
\titleformat{\subsubsection}{\normalsize\bfseries\color{blue!35!black}}{\thesubsubsection}{0.6em}{}
\titlespacing*{\subsubsection}{0pt}{1.2em}{0.5em}
\titleformat{\paragraph}[runin]{\normalsize\bfseries\color{blue!30!black}}{}{0pt}{}[]
\titlespacing*{\paragraph}{0pt}{1.0em}{0.6em}

\setlength{\parindent}{0pt}
\setlength{\parskip}{0.5em}
\allowdisplaybreaks

% ---------- Environments ----------
\theoremstyle{definition}
\newtheorem{definition}{Definition}[section]
\newtheorem{theorem}[definition]{Theorem}
\newtheorem{lemma}[definition]{Lemma}
\newtheorem{corollary}[definition]{Corollary}
\newtheorem{proposition}[definition]{Proposition}
\newtheorem{example}[definition]{Example}
\newtheorem*{remark}{Remark}
\newtheorem*{note}{Note}
\newtheorem*{claim}{Claim}
\newtheorem*{fact}{Fact}

% Visual framing of the theorem-like environments
\tcolorboxenvironment{definition}{enhanced,breakable,colback=blue!3,colframe=blue!3,
  borderline west={2pt}{0pt}{blue!55!black},arc=0pt,boxrule=0pt,left=6pt,right=4pt,top=2pt,bottom=2pt,before skip=10pt,after skip=10pt}
\tcolorboxenvironment{theorem}{enhanced,breakable,colback=red!2,colframe=red!2,
  borderline west={2pt}{0pt}{red!55!black},arc=0pt,boxrule=0pt,left=6pt,right=4pt,top=2pt,bottom=2pt,before skip=10pt,after skip=10pt}
\tcolorboxenvironment{lemma}{enhanced,breakable,colback=red!2,colframe=red!2,
  borderline west={2pt}{0pt}{red!55!black},arc=0pt,boxrule=0pt,left=6pt,right=4pt,top=2pt,bottom=2pt}
\tcolorboxenvironment{corollary}{enhanced,breakable,colback=red!2,colframe=red!2,
  borderline west={2pt}{0pt}{red!55!black},arc=0pt,boxrule=0pt,left=6pt,right=4pt,top=2pt,bottom=2pt}
\tcolorboxenvironment{proposition}{enhanced,breakable,colback=red!2,colframe=red!2,
  borderline west={2pt}{0pt}{red!55!black},arc=0pt,boxrule=0pt,left=6pt,right=4pt,top=2pt,bottom=2pt}
\tcolorboxenvironment{example}{enhanced,breakable,colback=white,colframe=white,
  borderline west={1.2pt}{0pt}{green!45!black},arc=0pt,boxrule=0pt,left=8pt,right=2pt,top=2pt,bottom=2pt,before skip=12pt,after skip=12pt}

% Emphasis for important parts
\newcommand{\term}[1]{\textbf{\color{blue!50!black}#1}}   % a term being defined
\newcommand{\key}[1]{\textbf{#1}}                         % key statement / takeaway phrase
% Syllabus week tag shown under a section heading
\newcommand{\syllabusweek}[1]{\par\vspace{-0.6em}{\small\color{gray}\textit{Syllabus: #1}}\par\vspace{0.4em}}
\setcounter{tocdepth}{2}

% Thin separator between logically distinct blocks inside a section
\newcommand{\sep}{\par\vspace{0.4em}\noindent{\color{gray!45}\rule{\linewidth}{0.4pt}}\par\vspace{0.2em}}
% Explanation of a step (replaces hand-written side notes): placed directly below/above the math
\newcommand{\explain}[1]{\par{\small\color{black!70}$\triangleright$\ \textit{#1}}\par}

% Added explanation (not in the original handwritten notes)
\newtcolorbox{insight}[1][Explanation]{enhanced,breakable,colback=yellow!4,colframe=orange!60!black,
  boxrule=0.4pt,arc=2pt,left=6pt,right=6pt,top=3pt,bottom=3pt,fonttitle=\small\bfseries,
  coltitle=orange!50!black,colbacktitle=yellow!12,title={#1},before skip=8pt,after skip=8pt}

% Highlighted box for results that are boxed in the handwritten notes
\newtcolorbox{keybox}[1][]{colback=blue!3,colframe=blue!50!black,boxrule=0.6pt,arc=2pt,#1}
% Side comment written in the margin / next to work in the notes
\newcommand{\sidenote}[1]{\quad\text{\small\textit{(#1)}}}
% Original page marker (for cross-checking with the scan)
\newcommand{\origpage}[1]{\leavevmode\marginnote{\raggedright\tiny\color{gray}p.~#1}}
% Illegible / uncertain reading
\newcommand{\unclear}[1]{{\color{red!70!black}[#1?]}}

% ---------- Math macros ----------
\newcommand{\R}{\mathbb{R}}
\newcommand{\C}{\mathbb{C}}
\newcommand{\N}{\mathbb{N}}
\newcommand{\Z}{\mathbb{Z}}
\newcommand{\Q}{\mathbb{Q}}
\newcommand{\F}{\mathcal{F}}
\newcommand{\V}{\mathcal{V}}
\newcommand{\W}{\mathcal{W}}
\newcommand{\Span}{\operatorname{sp}}
\newcommand{\Range}{\mathcal{R}}
\newcommand{\Null}{\mathcal{N}}
\DeclareMathOperator{\rank}{rank}
\DeclareMathOperator{\tr}{tr}
\DeclareMathOperator{\diag}{diag}
\DeclareMathOperator{\adj}{adj}
\newcommand{\Hm}{^{H}}            % Hermitian transpose
\newcommand{\T}{^{T}}             % transpose
\newcommand{\inner}[2]{\left\langle #1, #2 \right\rangle}
\newcommand{\norm}[1]{\left\lVert #1 \right\rVert}
\newcommand{\abs}[1]{\left\lvert #1 \right\rvert}
\newcommand{\0}{0_{\F}}
\newcommand{\1}{1_{\F}}
\newcommand{\bmat}[1]{\begin{bmatrix}#1\end{bmatrix}}
